\section{GTKP}
\subsection{Face up-reductions}
To any poset we can assign a topology compatible with the topological semantics.

\begin{definition}[Alexandrov topology]
    Let $(P,\leq)$ be a poset. We call the \textbf{Alexandrov} topology on $P$ the collection $\tau\subseteq \Pp P$ where $U\in\tau$ if and only if $U$ is upward-closed.
\end{definition}

We use it to define a weaker version of the up-reduction between polyhedra.
\begin{definition}[Face up-reduction]
    Let $P,Q$ be (compact) polyhedra and $K$ a triangulation of $Q$.

    A \textbf{face up-reduction} $P\fupred Q$ is a partial map $f:P\upred [K]$ such that
    \begin{enumerate}[(i)]
        \item the preimage of every open of $[K]$ (in the Alexandrov topology) is an open subpolyhedron of $P$;
        \item it is \textbf{open};
        \item it is \textbf{surjective}.
    \end{enumerate}
\end{definition}

We note the following:
\begin{proposition}
    $Q\upred[K]$ for any triangulation $K$.

    In particular, $Q\fupred Q$.
\end{proposition}
\begin{proof}
    Let $K$ be a triangulation of $Q$, and let $f$ be the map $q\mapsto \sigma_q$, where $\sigma_q$ is the unique face of $K$ such that $q\in\sigma_q$. Since every point of $Q$ lies in a unique such face, $f$ is trivially surjective.

    Note that the preimage of a face $\sigma$ is the corresponding face of the polyhedron, $f^{\gets}(\sigma)=|\sigma|$.

    If $U$ is an upward-closed subset of $[K]$, then $[K]\setminus U$ is downward-closed, and hence is a compact polyhedron. This means that $f^\gets([K]\setminus U)=Q\setminus f^{\gets}(U)$ is compact, and hence closed; therefore $f^\gets(U)$ is open and is a union of faces, i.e., an open subpolyhedron.

    This same argument shows that the map is open: if $P\subseteq Q$ is an open subpolyhedron, then $Q\setminus P$ is a compact polyhedron, so $[K]\setminus f(P)$ is downward-closed, so $f(P)$ is upward-closed, i.e., open.
\end{proof}
\begin{corollary}
    If there exists an up-reduction, i.e., an open partial simplicial map $P\upred Q$ whose domain is an open subpolyhedron of $P$, then $P\fupred Q$.
\end{corollary}


We may also make explicit the following fact:

\begin{lemma}
    Let $f:P\fupred Q$ be a face up-reduction, let $U\subseteq^o P$ be its domain, and let $L$ be the triangulation of $Q$ such that $f:U\upred [L]$.

    Then there exists a triangulation $K$ of $P$ such that
    \begin{itemize}
        \item $U$ is the union of relative interiors of faces of $K$, i.e. $U=\bigcup_{\sigma\in K_U}\RelInt\sigma$.
        \item $\tilde{f}:[K]\upred[L]$ defined as $\tilde{f}(\sigma_p)=f(p)$ is an up-reduction of posets.
    \end{itemize}
\end{lemma}
\begin{proof}
    Note that $[L]$ is a \textbf{finite} poset and that $\set{f^\gets(l):l\in L}$ is a collection of possibly open polyhedra in $P$. By the common triangulation lemma, we may extract a common triangulation $K$ of all of them.

    Now, recall that for each point $p\in P$ there exists a unique face $\sigma_p$ such that $p\in\RelInt(\sigma_p)$. This allows us to define $\tilde{f}$ as above. Uniqueness ensures that $\tilde{f}$ is well-defined, and openness of $f$ gives $f(\uparrow\sigma_p)=\uparrow f(p)$.
\end{proof}
\begin{lemma}[Up-reductions refute formulas]
    Let $P$ be a polyhedron, $S$ a poset, and $P\upred S$ an up-reduction.

    If $S\not\models \varphi$, then $P\not\models\varphi$ as well.
\end{lemma}
\begin{proof}
    The proof proceeds essentially as in the topological case in \cite{Gabelaia2001ModalDI}.
\end{proof}

\subsection{Goldblatt-Thomason for compact polyhedra}
\begin{theorem}[GTKP]
    A class of compact polyhedra $\C$ is modally definable if and only if it is closed under finite disjoint unions and face up-reductions, i.e.,
    \begin{enumerate}[(I)]
        \item if $P,Q\in \C$, then $P\sqcup Q\in \C$;
        \item if $P\in \C$ and $P\fupred X$, then $X\in \C$.
    \end{enumerate}
\end{theorem}
\begin{proof}[Proof ($\implies$)]
    Let $\C$ be a modally definable class, i.e., $\C=\set{P\in \Pol:P\models F}$ for some set of formulas $F$.

    Closure under disjoint union is trivial: if $P,Q\in \C$, then $P\sqcup Q\models F$.

    Suppose $P\in \C$ and $P\fupred X$, and suppose toward a contradiction that $X\notin \C$. By definition, this means that there exists a formula $\varphi\in F$ such that $X\not\models\varphi$.

    Spelling out the validity condition, there exists an evaluation $\nu:\Ll\to \Pp X$ and a point $x\in X$ such that $X,x\not\models\varphi$.

    By the common triangulation lemma, we may select a triangulation $S$ of $X$ such that $\nu(\varphi)$ is triangulated by $S$.

    Call $K=(V,\Sigma)$ the triangulation of $X$ realizing the face up-reduction $P\fupred X$. By the `nerves are dense' lemma, there exists a barycentric subdivision $K^{(n)}$ of $K$ subdividing $S$.

    We also know that we can extend any up-p-morphism $P\upred \Sigma$ to an up-p-morphism $P\upred \Sigma^{(k)}$, for any $k\in \N$.

    This gives an up-p-morphism $P\upred \Sigma^{(k)}\upred S$, and therefore $P\not\models\varphi$ --- a contradiction.
\end{proof}
\begin{proof}[Proof ($\impliedby$)]
    Let $\C$ be a class of polyhedra closed under (I) and (II); call $L:=\Log(\C)$ and $\C'=\Pol(L)$.

    Clearly $\C\subseteq \C'$. Let $X\in \C'$; we want to show that $X\in \C$.

    We know that $X\models L$; let $K=(V,\Sigma)\in \Tr(X)$ and call $v_i$ the vertices of $\Sigma$. We know that
    \[\bigsqcup_i (\uparrow v_i)\twoheadrightarrow \Sigma\]
    is a p-morphism.

    Consider $\chi_i:=\chi(\uparrow v_i)$, the Jankov-Fine formula for the (rooted) subposet $\uparrow v_i$.

    Since $(\uparrow v_i)\not\models \chi_i$, and each $(\uparrow v_i)$ is a generated subframe of $\Sigma$, we know that $\Sigma\not\models\chi_i$.

    In particular $X\upred\Sigma$, and thus $X\not\models\chi_i$ for all $i$. This means that none of the $\chi_i$ belong to $L$, so for each $\chi_i$ there exists some $P_i\in \C$ such that $P_i\not\models \chi_i$, and thus $P_i\upred(\uparrow v_i)$.

    It follows that $P:=\bigsqcup_i P_i$ admits an up-p-morphism to $\bigsqcup_i(\uparrow v_i)$; composing with $\bigsqcup_i(\uparrow v_i)\twoheadrightarrow \Sigma$ yields an up-reduction $P\upred \Sigma$, and hence a face up-reduction $P\fupred X$.

    By (I), $P\in \C$, and by (II), $X\in \C$.
\end{proof}

\begin{corollary}
    A class of polyhedra closed under finite disjoint unions and simplicial up-reductions is modally definable, where a simplicial up-reduction is an open simplicial partial map $f:P\upred Q$ whose domain is an open subpolyhedron of $P$.
\end{corollary}
\begin{proof}
    This follows from Corollary 3.1.4: if $f:P\upred Q$ is a simplicial up-reduction, then there exists a face up-reduction $P\fupred Q$ by composing with the face p-morphism onto any face poset.
\end{proof}

The theorem above can be strengthened by the following lemma.
\begin{lemma}[Any triangulation]
    Suppose $P\upred \Sigma$ is an up-reduction, for a triangulation $L=(V,\Sigma)$ of $Q$, and let $L'=(V',\Sigma')$ be a different triangulation of $Q$. Then there also exists an up-reduction $P\upred \Sigma'$.
\end{lemma}
\begin{proof}
    Let $f:P\upred\Sigma$ be an up-reduction with domain $U\subseteq P$. We extend $f$ to $f^{(1)}:P\upred\Sigma^{(1)}$ as follows.

    Take the triangulation $K$ of $P$ realizing $f$ as a simplicial map $\tilde{f}$ (Lemma 3.1.5), and use the subdivision lemma (Lemma 1.3.4) to extract an up-reduction $K'\upred L^{(1)}$.

    Now, repeatedly applying Lemma 1.3.3, we may iterate this $n$ times so that $L^{(n)}$ is a subdivision of $L'$. This gives a map $h: K^{\prime \cdots\prime }\upred L^{(n)}\twoheadrightarrow L'$.

    Since $K^{\prime \cdots\prime }$ is a triangulation of $P$, we can extract an up-reduction $g:P\upred L'$ by setting it to be locally constant on the relative interior of each face of $P$.
\end{proof}

In other words, if a face up-reduction exists to one triangulation, then one exists to any triangulation, even with the same domain.
